\documentclass[11pt]{article}
\usepackage[T1]{fontenc}
\usepackage[utf8]{inputenc}
\usepackage{lmodern,amsmath,amssymb,amsthm}
\usepackage[margin=1in]{geometry}
\usepackage[hidelinks]{hyperref}
\newtheorem{lemma}{Lemma}\newtheorem{theorem}[lemma]{Theorem}
\newcommand{\R}{\mathbb R}\newcommand{\Z}{\mathbb Z}
\newcommand{\C}{\mathbb C}\newcommand{\E}{\mathbb E}
\newcommand{\ee}{\mathrm e}
\title{The polynomial GI(2) base from the published bracket theorem}
\author{Author proposal: /root/astra\_v14\_baseline\_theta}
\date{7 November 2026}
\begin{document}\maketitle

\noindent\textbf{Scope.} This is an author proposal awaiting independent
review. It derives the exact polynomial interval base from the forward
bracket-quadratic theorem and Bohr regularity in Green--Tao,
\emph{An inverse theorem for the Gowers $U^3(G)$ norm}, Proceedings of
the Edinburgh Mathematical Society 51 (2008), 73--153. It replaces the
needed use of the false quadratic-extension assertion by an explicit
finite-box proof. It also constructs the bracket tests directly on
connected polynomial models, so neither Theorem 12.9 nor its erroneous
skew-shift iterate parameters are a theorem premise. The forward
Bohr/bracket inverse theorem and its unchanged proof inputs remain
published inputs, not newly proved by this unit. The complete GI theorem
and broader source-risk review remain separate.

\section{Exact target and published input}

The required base is this statement. For every $\delta>0$ there are
finitely many fixed rational filtered nilmanifolds of degree at most two,
with connected simply connected groups, and constants $L<\infty,c>0$
such that every complex $1$-bounded $f:[N]\to\C$ with
$\|f\|_{U^3[N]}\geq\delta$ correlates by at least $c$ with
$F(g(n)\Gamma)$ on one of them. The polynomial $g$ is adapted,
$\|F\|_\infty\leq1$, and $\operatorname{Lip}(F)\leq L$.
All geometry and constants are fixed before $N$ and $f$.
The interval norm is zero extension to a sufficiently large cyclic
group divided by the norm of the interval indicator.

The published forward input is Theorem 10.11, pp.~124--125, especially
(10.7) and (10.9). For each fixed sufficiently small $\eta>0$, and
each prime $P>2$, a function $u:\Z/P\Z\to\C$ with
$|u|\leq1$ and $\|u\|_{U^3}\geq\eta$ yields a frequency set
$S$ of size at most $D(\eta)$, a proper progression $Q$ containing
a Bohr set $B(S,\rho_0)$ with $\rho_0\geq\rho_*(\eta)>0$,
and one bracket quadratic
\[
 \phi(n)=c+\sum_j a_j\{\alpha_jn\}
                  +\sum_{j,k}a_{jk}\{\alpha_jn\}\{\alpha_kn\}
       \pmod1 .
\]
Here $\alpha_j=\xi_j/P$ represents a frequency in $S$,
$\{t\}\in[-1/2,1/2)$, and coefficients are unrestricted real numbers.
For every nonempty $A\subseteq Q$ there is a shift $h_A$ such that
\[
 \left|\E_{n\in A}u(n+h_A)\ee(-\phi(n))\right|
 \geq b_*(\eta)>0.
\]
The actual powers and exponential constants printed in 10.11 are
stronger than needed; only the stated positive/finiteness dependence
on $\eta$ is consumed. Its proof uses the forward Theorem 2.7,
Corollary 10.6 and Proposition 10.10. The affected special case in the
latter is supplied next.

\section{The exact specialization replacing quadratic extension}

The final PDF's Lemma 10.7, pp.~122--123, invokes quadratic extension
from a subgroup. The author notice of December 2023 identifies
Proposition 3.2, p.~88, as false. The later use needed here is
Proposition 10.10, where the progression is proper and $H=\{0\}$.
The notice numbers these items 10.6 and 10.8 in its preprint version;
we bind to the final PDF numbering and argument.

\begin{lemma}\label{lem:box}
On a finite box $B=\prod_{j=1}^d\{0,\ldots,L_j-1\}$, any
$v:B\to\R/\Z$ whose third additive differences vanish whenever
the whole cube lies in $B$ has the form
\[
 v(l)=c+\sum_j b_jl_j+\sum_j d_j\binom{l_j}{2}
                       +\sum_{j<k}e_{jk}l_jl_k.
\]
Consequently a locally quadratic phase on a proper progression with
trivial subgroup has the polynomial coordinate representation needed
in Proposition 10.10, without extension to the ambient group.
\end{lemma}
\begin{proof}
Coordinates of length one can be omitted. Each mixed second difference
is constant on its rectangular domain: every coordinate difference of
it is a vanishing third difference, and the domain is connected by
coordinate steps. The same holds for each available diagonal second
difference when $L_j\geq3$. Choose their constants as $e_{jk}$ and
$d_j$, setting $d_j=0$ when $L_j=2$. Subtract these quadratic terms.
Every remaining first difference is constant because its differences
in every available coordinate vanish. When $L_j=2$ its domain has only
one slice in direction $j$, so mixed differences suffice. Subtract the
resulting linear function; the remainder is constant by connectedness.
The coordinate map of a proper progression is injective and pulls
back the local quadratic identities to the box. This proves the
claim. Division by two in $\R/\Z$, if one wants the symmetric
coefficient convention of (10.5), is possible; no uniqueness is used.
\end{proof}

In Proposition 10.10 the independent vectors supplied by Corollary 10.6
express each progression coordinate as a real linear combination of
$\{\xi n/P\}$ on that progression. Substituting the lemma's polynomial
gives precisely its bracket quadratic. The later localization and
shift in the proof of 10.11 use this specialization, not the extension
of a phase from a nontrivial subgroup. The remaining published forward
inverse and Bohr-set applications are retained as explicit inputs.
This does not claim to prove a corrected general version of 3.2.

\section{A smooth periodic bracket test with positive correlation}

Use Lemmas 8.1--8.2 to choose a regular Bohr set
$B=B(S,\rho)\subseteq Q$, where $\rho$ has a positive lower bound
depending only on $\eta$ and may be taken at most $1/8$.
Its density in $\Z/P\Z$ is at least $b_0(\eta)>0$.
Apply the displayed (10.9) input with $A=B$ and write its shift as $h$.

For $d=|S|>0$, choose a fixed $0<\epsilon<1/(100d)$ sufficiently
small in terms of $b_*(\eta)$ and $d$. Let $\chi:\R/\Z\to[0,1]$
be one on $|x|\leq\rho(1-\epsilon)$, zero on
$|x|\geq\rho(1+\epsilon)$, and linear on the intervening strips.
Its Lipschitz bound is at most $(2\epsilon\rho)^{-1}$, with an
inessential larger fixed bound permitted. In particular it vanishes
on a fixed neighbourhood of the fractional-part seam at $1/2$.
Put
\[
 W(n)=\prod_{j=1}^d\chi(\alpha_j n)^{1+4d},\qquad
 \Psi(n)=W(n)\ee(-\phi(n)).
\]
The difference $W-1_B$ vanishes outside the regularity boundary shell
$B(S,\rho(1+\epsilon))\setminus B(S,\rho(1-\epsilon))$ and
is bounded by one. Bohr regularity bounds the shell by
$C d\epsilon |B|$ with an absolute fixed $C$. Choose $\epsilon$
so that $Cd\epsilon\leq b_*(\eta)/2$. It follows that
\begin{equation}\label{eq:bias}
 \left|\E_{n\in\Z/P\Z}u(n+h)\Psi(n)\right|
 \geq \tfrac12 b_*(\eta)b_0(\eta)=:c_0(\eta)>0.
\end{equation}
When $S$ is empty the phase is constant and this conclusion already
uses a constant test; that case needs no cutoff construction.

Every factor in $\Psi$ is periodic with period $P$, because
$P\alpha_j$ is integral. We now represent this same function for
all integers, not merely signed representatives of $\Z/P\Z$.

\section{Connected polynomial models for every bracket factor}

Write the coefficients below as $a=q+s$, with $q\in\Z$ and
$|s|\leq1/2$. A linear bracket factor is
\[
 \chi(\alpha n)\ee(a\{\alpha n\})
 =\chi(\alpha n)\ee(s\{\alpha n\})\ee(q\alpha n).
\]
It is the test $F(x,y)=\chi(x)\ee(sx+y)$ along the linear
polynomial $(\alpha n,q\alpha n)$ on $\R^2/\Z^2$.
The seam is harmless because $\chi$ vanishes there. This test is
$1$-bounded and has Lipschitz bound $O(1+\operatorname{Lip}\chi)$,
uniformly in $a,\alpha$.

A quadratic bracket factor satisfies the exact identity
\begin{align*}
 &\chi(\alpha n)^2\chi(\gamma n)^2
                    \ee(a\{\alpha n\}\{\gamma n\})\\
 &=\bigl[\chi(\alpha n)\chi(\gamma n)
             \ee(s\{\alpha n\}\{\gamma n\})\bigr]
   \ee(q\alpha\gamma n^2)
   \bigl[\chi(\gamma n)\ee(-q\alpha n[\gamma n])\bigr]
   \bigl[\chi(\alpha n)\ee(-q\gamma n[\alpha n])\bigr].
\end{align*}
Indeed the omitted term in expanding the integer part is
$q[\alpha n][\gamma n]\in\Z$, where $[t]=t-\{t\}$.
The first bracket is a test on a two-torus with uniformly bounded
Lipschitz constant. The pure quadratic phase is a test $\ee(z)$
along $z=q\alpha\gamma n^2$ on $\R/\Z$, using the degree-two
filtration $G_0=G_1=G_2=\R$, $G_3=0$.

For the remaining bracket use the connected simply connected Heisenberg
group with coordinates $(u,v,w)$ and law
\[
 (u,v,w)(u',v',w')=(u+u',v+v',w+w'+uv').
\]
Its lattice is $\Z^3$ in these coordinates; its degree-two filtration
has $G_0=G_1=G$ and $G_2=\{(0,0,w)\}$.
On a right lattice quotient, reduce $v$ first, then $u$. The resulting
coordinates are
\[
 (\{u\},\{v\},w-u[v]\pmod1).
\]
Define a test on this fundamental cube by
$F_H(\{u\},\{v\},z)=\chi(\{v\})\ee(z)$.
It is well-defined and globally Lipschitz. Across the $u$ faces the
central coordinate does not change; the central faces are periodic;
across the $v$ faces the cutoff is identically zero in a neighbourhood.
Within the remaining finitely many fixed charts its Lipschitz constant
is bounded by a fixed geometric multiple of
$1+\operatorname{Lip}\chi$.

Evaluate this test on the polynomial
$p(n)=(\beta n,\alpha n,0)$. The exact value is
\[
 F_H(p(n)\Gamma)=\chi(\alpha n)\ee(-\beta n[\alpha n]).
\]
For completeness its left first difference is
$(\beta h,\alpha h,-\alpha\beta hn)$, its next difference
is $(0,0,-\alpha\beta hk)$, and its third difference is the
identity. These satisfy the degree-two filtration for every integer
shift, so this is an adapted polynomial without any appeal to an
incorrect formula for linear skew iterates. Take $\beta=q\gamma$
and interchange $\alpha,\gamma$ for the two remaining factors.

Apply the linear construction to each coefficient of $-\phi$ and
the quadratic construction to all ordered pairs, inserting zero
coefficients if needed. Each frequency receives one cutoff from the
linear factor and $4d$ from the quadratic factors, exactly the power
in $W$. The constant phase is harmless. Products give a fixed
connected rational polynomial model of degree at most two, of dimension
at most $2d+9d^2$. The product test is $1$-bounded and has a finite
Lipschitz bound depending only on $d,\rho,\epsilon$. Arbitrarily
large coefficients occur only in polynomial orbit parameters.
There are finitely many geometric product types for $d\leq D(\eta)$.

Thus the model represents the periodic function $\Psi(n)$ exactly
for all $n\in\Z$. Individual polynomial factors need not be
$P$-periodic; their product is periodic by the exact bracket identity.
That distinction is used in the next change of variables.

\section{The normalized interval transfer}

Fix an integer $C>16$ and choose a prime $P$ with $CN<P<2CN$.
Zero-extend the given interval function to $u$ on $\Z/P\Z$.
The interval-indicator Gowers norm is at least its mean $N/P$, so
$\|u\|_{U^3}\geq\delta/(2C)$. Choose once and for all a positive
$\eta$ at most this bound and within the smallness range of the
published forward theorem. All preceding data now depend only on
$\delta$.

Using the exact periodicity of $\Psi$ in \eqref{eq:bias}, change
variables to obtain
\[
 \left|\E_{n\in\Z/P\Z}u(n)\Psi(n-h)\right|\geq c_0(\eta).
\]
The zero extension gives
\[
 \left|\E_{n\in[N]}f(n)\Psi(n-h)\right|
 \geq (P/N)c_0(\eta)>0.
\]
The translated sequence is represented by the same test on the same
fixed model, with polynomial orbit $g(n-h)$. Translation preserves
adapted polynomiality, either by the derivative criterion or B.5.
No signed-representative split, discontinuous interval indicator or
additional $N$-dependent geometry is introduced. Norm, Lipschitz
and correlation bounds have precisely the required fixed dependence.
Bertrand's postulate applies to the integer $CN>1$, so the prime
choice also covers every positive $N$.

\section{Source comparison and bounded related-use findings}

The unpublished author notice identifies the false general quadratic
extension, but is used here only to locate the affected proof. The
needed trivial-subgroup specialization is proved in Section 2.
Its published source locator is Proposition 3.2 on p.~88, correcting
the p.~89 metadata typo in assembly001's RELATED-USES record.

The actual publisher image of Example 12.3 on p.~136 gives
$T(x,y)=(x+\alpha,y+\beta+mx)$ but prints an iterate omitting
$mnx$ and using $mn(n+1)\alpha/2$. Direct iteration gives
\[
 T^n(x,y)=(x+n\alpha,
 y+n\beta+mnx+m\alpha\tbinom n2).
\]
Already $n=1$ distinguishes these formulas. The later pure quadratic
choice in Lemma 12.11 inherits the printed iterate convention.
Rather than silently change its parameters, Section 4 constructs
all needed factors directly on connected polynomial models.
This preserves the qualitative base theorem's strength and avoids
those printed iterates entirely.

The cutoff's derivative bound also depends on its transition width.
This proof retains the finite factors $(\epsilon\rho)^{-1}$ instead
of importing the universal numerical bounds stated for those tests
in Lemma 12.11. The required application needs only finite bounds
after $\eta$ is fixed.

The bounded audit scope is the forward chain
$2.7\to10.6\to10.10\to10.11$, the affected trivial-subgroup use
of 10.7, Bohr regularity 8.1--8.2, and the elementary construction
12.2--12.11 as consumed by the interval base. The chain notation
records dependencies to inspect, not acceptance of all their proofs.
The converse half of 2.7, the finite-field-only extension use and
unrelated applications are not premises of this unit. The old
assembly001/base text remains frozen. This successor changes its
base supplier from 12.9 to the precise bracket theorem plus local
construction; its separate converse argument is unchanged.

\end{document}
